Tato diplomová práce se zabývá algoritmy strojového učení pro adaptaci autonomního agenta na chování uživatele ve 3D akční hře. Hlavním cílem je navrhnout jednoduchého autonomního agenta, který se pomocí strojového učení a~změny vybavení dokáže přizpůsobit protivníkovi a~porazit ho.

První kapitola obsahuje dvě podkapitoly. První se věnuje vybraným technikám a~algoritmům strojového učení, jako jsou Proximal Policy Optimization (PPO), Deep Q-Network (DQN), a~techniky Self-Play a~Population Based Training (PBT). Druhá podkapitola popisuje existující herní agenty, konkrétně AlphaStar, OpenAI Five a~Voyager, včetně přehledu jejich fungování a~způsobu implementace.

Druhá kapitola se zaměřuje na použité technologie a~návrh samotné hry. První podkapitola podrobně popisuje využití Unreal Engine, Visual Studia a~pluginu Nevarok ML, který rozšiřuje Unreal Engine o~funkcionalitu pro trénink a~aplikaci modelů strojového učení. Druhá podkapitola se věnuje návrhu herního prostředí a~autonomního agenta.

Třetí kapitola popisuje implementaci hry a~agenta. První část se zaměřuje na herní logiku, která poskytuje prostředí pro aplikování agenta. Druhá část se věnuje implementaci samotné logiky agenta, pokusům o~vylepšení jeho logiky a~porovnání různých verzí.

Závěr shrnuje úvahy autora a~dosažené výsledky — zda se agent adaptuje podle očekávání, jaký vliv mají jednotlivá vylepšení a~jak si různé verze modelu vedou ve vzájemném porovnání.